\section{中心极限定理}

本节介绍中心极限定理。
中心极限定理描述了随着试验次数的增加，所有试验结果的值的累计趋于正态分布这一现象。

本节要点：
\begin{itemize}
    \item 理解列维—林德伯格中心极限定理。
\end{itemize}

%============================================================
\subsection{列维—林德伯格中心极限定理}

\begin{theorem}[列维—林德伯格中心极限定理]
设$X_1,X_2,\cdots ,X_n,\cdots $为独立同分布的随机变量，若它们都有相同的数学期望和方差，分别记为$\mu ,\sigma ^2$，有：
\[
\underset{n\rightarrow \infty}{\lim}P\left\{ \frac{X_1+X_2+\cdots +X_n-n\mu}{\sqrt{n}\sigma}\leqslant x \right\} =\varPhi \left( x \right) =\frac{1}{\sqrt{2\pi}}\int_{-\infty}^x{e^{-\frac{t^2}{2}}dt}
\]
即当$n\rightarrow \infty $时，$\frac{X_1+X_2+\cdots +X_n-n\mu}{\sqrt{n}\sigma}$服从标准分布。
\end{theorem}

列维—林德伯格中心极限定理告诉我们，若有一系列独立同分布的随机变量，倘若将它们相加后再做标准化：
\[
\begin{cases}
	Y=X_1+X_2+\cdots +X_n\\
	EY=n\mu\\
	DY=n\sigma ^2\\
\end{cases}\Rightarrow Y^*=\frac{Y-EY}{\sqrt{DY}}=\frac{Y-n\mu}{\sqrt{n}\sigma}
\]
则当$n\rightarrow \infty $时$Y^*\sim N\left( 0,1 \right) $，也即$Y\sim N\left( n\mu ,n\sigma ^2 \right) $。
实际中，一般$n\geqslant 30$时，我们就可以认为$X_1,X_2,\cdots ,X_n$的叠加服从正态分布，即：
\[
\sum_{i=1}^n{X_i}\sim N\left( n\mu ,n\sigma ^2 \right)
\]

中心极限定理更重要的实际意义在于告诉我们，$X_1,X_2,\cdots ,X_n$只需要独立同分布就行，具体是什么分布并不重要，可以是任意的、未知的、或者不可表达式化的，甚至它们之间允许有略微的非独立，只要$n$够大，它们的叠加就可以看成正态分布。

%============================================================
\subsection{中心极限定理的说明}

除了列维—林德伯格中心极限定理，还有棣莫弗—拉普拉斯中心极限定理等。
中心极限定理是当$n$足够大时叠加可以看成正态分布的一系列定理的统称。
中心极限定理依然要求“独立试验”。

%============================================================
\subsection{例}

\begin{example}
设有200台独立工作的设备，每台设备的故障概率为1/200，如果出现故障，只需一名维修工即可修复，试分析需要多少维修工才能使出现故障无法及时修理的概率控制在0.03以内。
\end{example}

解：

分析，由于一台故障设备只需一名维修工修复，所以问题就成为这200台设备同时有几台及以上同时发生故障的概率$<3\%$，也即这200台设备同时有几台及以下同时发生故障的概率$\geqslant 97\%$。

1. 设备故障可以视作0-1分布且$p=1/200$，令记$X_i=1$表示第$i$台设备出现故障，则$X_i$服从0-1分布，即
\begin{align*}
&X_i\sim B\left( 1,p \right) \\
&p=\frac{1}{200} \\
&\mu =EX_i=p=\frac{1}{200} \\
&\sigma ^2=DX_i=p\left( 1-p \right) =\frac{199}{200\cdot 200}
\end{align*}

2. 令$Y=\sum{X_i}$，即表示200台设备中同时故障的台数，$Y$也是随机变量，于是根据列维—林德伯格中心极限定理，$Y$服从正态分布$N\left( n\mu ,n\sigma ^2 \right) $：
\[
Y\sim N\left( n\mu ,n\sigma ^2 \right) \overset{n=200}{=}N\left( 1,\frac{199}{200} \right)
\]

3. 根据正态分布和标准正态分布的关系，$y$及$y$台以下同时故障的概率：
\[
P\left\{ Y\leqslant y \right\} =\varPhi \left( \frac{y-\mu}{\sigma} \right) =\varPhi \left( \frac{y-1}{\sqrt{199/200}} \right) \geqslant 97\%
\]

4. 查表得$\varPhi \left( 1.88 \right) \approx 0.97$，即$P\left\{ Y\leqslant 2.88 \right\} =0.97$，说明3台及3台以下设备同时故障的概率占了97\%以上，因此只要配备3个维修工即可。

解二：
也可以用泊松分布，令$X$为同时故障的台数，服从泊松分布：
\[
P\left\{ X=x \right\} =\frac{\lambda ^x}{x!}e^{-\lambda},\quad \lambda =200\cdot \frac{1}{200}=1
\]
$n$台及$n$以下同时故障的概率：
\[
P\left\{ X\leqslant n \right\} =\sum_{x=0}^n{\frac{1}{x!}e^{-1}}
\]
计算$n=1,2,3,4$得：
\begin{align*}
&P\left\{ X\leqslant 1 \right\} =0.736 \\
&P\left\{ X\leqslant 2 \right\} =0.920 \\
&P\left\{ X\leqslant 3 \right\} =0.981 \\
&P\left\{ X\leqslant 4 \right\} =0.996
\end{align*}
3台及以上同时故障的概率为$1-0.981=1.9\%<3\%$，后同，略。

解一和解二的分布如下图，可见当$n$足够大时，两者分布趋于一致。
\begin{figure}[h]
\centering
\includegraphics[height=5cm]{5.1.png}
\end{figure}

~

\begin{example}
同上，现设4名维修工，有两个方案，(A) 4名维修工各负责50台设备，(B) 4名维修工分2组，每2名一组，每组负责100台设备，试分析哪个方案较为合理。
\end{example}

解：

每台故障依然如上：
\begin{align*}
&X_i\sim B\left( 1,p \right) \\
&p=\frac{1}{200} \\
&\mu =EX_i=p=\frac{1}{200} \\
&\sigma ^2=DX_i=p\left( 1-p \right) =\frac{199}{200\cdot 200}
\end{align*}

(A) 当每人负责50台设备时，相当于将200台设备分4组，令每组故障数为随机变量$Y_i,i=1,2,3,4$，则有：
\begin{align*}
&Y_i\sim N\left( n\mu ,n\sigma ^2 \right) \overset{n=50}{=}N\left( \frac{50}{200},\frac{50\cdot 199}{200\cdot 200} \right) \\
&P\left\{ Y_i\leqslant y \right\} =\varPhi \left( \frac{y-\mu}{\sigma} \right) =\varPhi \left( \frac{y-\frac{50}{200}}{\sqrt{\frac{50\cdot 199}{200\cdot 200}}} \right) =\varPhi \left( \frac{y-1/4}{\sqrt{199/800}} \right)
\end{align*}
要使设备故障时都能及时得到维修，则
\begin{align*}
P&=P\left\{ Y_1\leqslant 1,Y_2\leqslant 1,Y_3\leqslant 1,Y_4\leqslant 1 \right\} \\
&=P\left\{ Y_1\leqslant 1 \right\} \cdot P\left\{ Y_2\leqslant 1 \right\} \cdot P\left\{ Y_3\leqslant 1 \right\} \cdot P\left\{ Y_4\leqslant 1 \right\} \\
&=\varPhi \left( \frac{1-1/4}{\sqrt{199/800}} \right) ^4=\varPhi \left( 1.5038 \right) ^4=0.9332^4 \\
&=0.758
\end{align*}

(B) 当每2人负责100台设备时，相当于将200台设备分2组，令每组故障数为随机变量$Y_i,i=1,2$，则有：
\begin{align*}
&Y_i\sim N\left( n\mu ,n\sigma ^2 \right) \overset{n=100}{=}N\left( \frac{100}{200},\frac{100\cdot 199}{200\cdot 200} \right) \\
&P\left\{ Y_i\leqslant y \right\} =\varPhi \left( \frac{y-\mu}{\sigma} \right) =\varPhi \left( \frac{y-\frac{100}{200}}{\sqrt{\frac{100\cdot 199}{200\cdot 200}}} \right) =\varPhi \left( \frac{y-1/2}{\sqrt{199/400}} \right)
\end{align*}
要使设备故障时都能及时得到维修，则
\begin{align*}
P&=P\left\{ Y_1\leqslant 2,Y_2\leqslant 2 \right\} \\
&=P\left\{ Y_1\leqslant 2 \right\} \cdot P\left\{ Y_2\leqslant 2 \right\} \\
&=\varPhi \left( \frac{2-1/2}{\sqrt{199/400}} \right) ^2=\varPhi \left( 2.1266 \right) ^2=0.983^2 \\
&=0.966
\end{align*}

可见，方案B故障及时维修的概率更大，更合理。




